\chapterhead{45}{ORR-6}{Risk Reporting}{1}{2}

\lohead{45}{a}{Identify roles and responsibilities of different organizational
committees, and explain how risk reports should be developed for each committee or
business function.}

\orsub{Organizational committees}

\noindent\begin{minipage}{\textwidth}
\renewcommand{\arraystretch}{1.3}
\begin{tabularx}{\textwidth}{@{}L{0.26\textwidth}YY@{}}
\rowcolor{navy} \thd{Committee or function} & \thd{Main role} & \thd{What it receives or reports}\\
{\tabfont \textbf{1) Risk Committee}} &
{\tabfont Oversee the firm's risk management framework.} &
{\tabfont Receives summary reports on risk exposures, trends, and emerging risks,
and provides directives based on the information received.}\\
\rowcolor{paleblue}
{\tabfont \textbf{2) Audit Committee (Board)}} &
{\tabfont Oversees internal controls.} &
{\tabfont Receives reports from internal audit on risk and control
effectiveness.}\\
{\tabfont \textbf{3) Executive Committee (ExCo)}} &
{\tabfont Steers the board, oversees risk management.} &
{\tabfont Receives reports on risk exposures, trends, events, action plans,
culture, and remediation efforts.}\\
\rowcolor{paleblue}
{\tabfont \textbf{4) Central Risk Function}} &
{\tabfont Collects and aggregates risk data.} &
{\tabfont Reports to the risk committee and business line managers. Collects data
on risk exposures, action plan statuses, controls, risk profile changes, and risk
events for comprehensive decision support.}\\
{\tabfont \textbf{5) Business Line Managers and Risk Champions}} &
{\tabfont Front-line monitoring of operational risk data and action plan
progress.} &
{\tabfont Focus on KRIs, action plan progress, types, and impacts of operational
risks. Monitor and benchmark own risks against firmwide averages or other lines
and departments.}\\
\end{tabularx}
\end{minipage}
\orfigcap{Who reports what, to whom.}

% ---------------------------------------------------------------------
\lohead{45}{b}{Describe components of operational risk reports and explain best
practices in operational risk reporting.}

\orsub{Operational risk reporting components}

\orsubb{The reporting cake (3 tiers)}

\begin{center}
\begin{tikzpicture}[font=\fontsize{8.2}{10.2}\selectfont,
  lbl/.style={text width=4.4cm,align=right,text=navy,
              font=\sffamily\bfseries\fontsize{8.2}{10}\selectfont},
  slab/.style={rounded corners=2.5pt,inner sep=7pt,align=center,text=white,
               font=\sffamily\bfseries\fontsize{8.6}{10.4}\selectfont}]

\node[slab,fill=navy,   text width=4.6cm] (s3) at (2.6,0)    {All you need to know to decide};
\node[slab,fill=teal,   text width=6.4cm] (s2) at (3.5,-1.35) {All you need to know to act};
\node[slab,fill=forest, text width=8.2cm] (s1) at (4.4,-2.70) {All you need to know to monitor};

\node[lbl,anchor=east] (l3) at (-2.3,0)    {Tier 3\\[1pt]{\mdseries\fontsize{7.6}{9.4}\selectfont Executive Committee, Board Risk Committee}};
\node[lbl,anchor=east] (l2) at (-2.3,-1.35) {Tier 2\\[1pt]{\mdseries\fontsize{7.6}{9.4}\selectfont Operational Risk Function, Operational Risk Committee}};
\node[lbl,anchor=east] (l1) at (-2.3,-2.70) {Tier 1\\[1pt]{\mdseries\fontsize{7.6}{9.4}\selectfont Risk Champions, Business-Line Managers}};

\foreach \a/\b in {l3/s3,l2/s2,l1/s1}
  \draw[palenavy,line width=1pt] (\a.east) -- (\b.west);
\end{tikzpicture}
\end{center}
\orfigcap{The reporting cake: detail narrows as reports travel upward.}

\orsub{Comprehensive ORM report components}

\begin{enumerate}
\item[\textcolor{rust}{\sffamily\bfseries 1)}] \textbf{Prioritised risks and
outlook.} Identifies critical risks such as technology failures, cyberattacks, and
regulatory breaches, ranking them by inherent or residual risk and providing an
outlook for each.
\item[\textcolor{rust}{\sffamily\bfseries 2)}] \textbf{Heatmap and risk register.}
Utilises risk and control self-assessment (RCSA) to list operational risks and
controls, categorising risks into a heatmap based on likelihood and severity,
using colours to indicate urgency levels.
\item[\textcolor{rust}{\sffamily\bfseries 3)}] \textbf{Risk appetite metrics.}
Presented to the board to ensure alignment with the firm's risk appetite. These
metrics track compliance with risk limits through KRIs, which include details such
as type, name, threshold, actual values, and scores.
\item[\textcolor{rust}{\sffamily\bfseries 4)}] \textbf{KRIs and issue monitoring.}
Used for forward-looking analyses and detailed risk assessments. There is also a
KRI dashboard which shows key metrics, trends, and scores to efficiently leverage
existing data. Examples of metric sources: \emph{human resources} (attendance at
compliance training, trained personnel percentage) and \emph{risk control}
(completion rates of risk management plans, timeliness of operational incident
reports).
\item[\textcolor{rust}{\sffamily\bfseries 5)}] \textbf{Risk events and near
misses.} Risk event reports include details on the events (number and size),
frequency and severity on a per-period basis, trends over time (monthly,
quarterly, or annually), and specific narratives concerning larger incidents that
extend above predefined materiality thresholds. Near misses are critical to report
and are evaluated based on the materiality of potential impact avoided. The
lessons learned are tremendously valuable to an entity, without the incurrence of
actual losses.
\item[\textcolor{rust}{\sffamily\bfseries 6)}] \textbf{Action plans and
remediation.} An action plan is designed to improve the control environment
through the use of risk mitigation programs.
  \begin{lettered}
  \item \textbf{Corrective plans:} react to large incidents or gaps in the ORM
  system.
  \item \textbf{Detective controls:} identify potential incidents before they
  escalate.
  \item \textbf{Preventive plans:} proactively manage risks to prevent major
  issues.
  \end{lettered}
\item[\textcolor{rust}{\sffamily\bfseries 7)}] \textbf{Emerging risks and horizon
scanning.} Identifies upcoming risks and trends for board risk committee reviews.
The focus here is on regulatory changes, compliance, and market volatility.
\item[\textcolor{rust}{\sffamily\bfseries 8)}] \textbf{Risk software.} Governance,
risk, and compliance (GRC) systems are risk software applications that help
organisations collect and report on operational risk data.
\end{enumerate}

% ---------------------------------------------------------------------
\lohead{45}{c}{Describe challenges to reporting operational risks, including
characteristics of operational loss data, and explain ways to overcome these
challenges.}

\orsub{Operational risk reporting challenges}

\begin{enumerate}
\item[\textcolor{rust}{\sffamily\bfseries 1.}] \textbf{Operational risk event data
asymmetry.} Operational losses are often characterised by a small number of
high-severity, low-frequency events. Therefore the focus in risk management should
prioritise the mitigation and prevention of significant incidents rather than
minor, frequent ones.
\item[\textcolor{rust}{\sffamily\bfseries 2.}] \textbf{Large risk event
escalation.} Upper management must closely assess and respond to significant risk
events and near misses that exceed the organisation's risk tolerance.
\item[\textcolor{rust}{\sffamily\bfseries 3.}] \textbf{Small, frequent losses.}
Most operational risk incidents are small but occur frequently. While they may not
significantly impact the budget, regular analysis is necessary to identify any
structural flaws or control breaches that require corrective action.
\item[\textcolor{rust}{\sffamily\bfseries 4.}] \textbf{Operational loss
benchmarking.} Benchmarking operational risk losses helps in capital allocation
and budgeting decisions. Benchmarks can be based on regulatory capital (measured
in basis points of capital consumed) or related to financial metrics like gross
income.
\item[\textcolor{rust}{\sffamily\bfseries 5.}] \textbf{Operational risk
distributions.} Operational risk data typically do not fit normal distributions,
making median and quartiles more useful than averages for analysis.
\item[\textcolor{rust}{\sffamily\bfseries 6.}] \textbf{Concentrations, outliers,
and scenarios.} Outliers are key for understanding patterns and data
distributions. Past patterns, distribution concentrations, and distances between
observations all merit significant analysis. Assessing various risk scenarios
(high, medium, low) is essential, with recent emphasis on modelling for
climate-related risks due to their increasing importance.
\item[\textcolor{rust}{\sffamily\bfseries 7.}] \textbf{Qualitative risk data
aggregation.}
  \begin{lettered}
  \item \textbf{Conversion and addition:} transform qualitative data into monetary
  values for clear financial representation.
  \item \textbf{Categorisation:} uses colour coding (red, yellow, green) to
  visually represent risk levels in reports.
  \item \textbf{Worst-case reporting:} applies a conservative approach, marking an
  entire group as high-risk if any single risk is high.
  \end{lettered}
\item[\textcolor{rust}{\sffamily\bfseries 8.}] \textbf{Combined assurance.}
  \begin{lettered}
  \item The main purpose is to coordinate assurance activities across the
  organisation to inform leadership about governance and risk management.
  \item \textbf{Ownership:} managed by the ORM function using a combined assurance
  map to report risk statuses (green, yellow, red).
  \item \textbf{Three lines of defense:} involves direct assessments, oversight
  functions, and periodic audits to ensure effective risk management.
  \end{lettered}
\end{enumerate}

% ---------------------------------------------------------------------
\lohead{45}{d}{Explain best practices for reporting risk exposures to regulators
and external stakeholders.}

\orsub{External reporting best practices}

The Basel regulatory framework (Pillar 3) addresses public disclosure of financial
and operational risk information. The three types of operational risk disclosure
requirements include the following:

\begin{enumerate}
\item[\textcolor{rust}{\sffamily\bfseries 1.}] Qualitative information.
\item[\textcolor{rust}{\sffamily\bfseries 2.}] Historical losses, which includes
10 years of aggregate operational losses.
\item[\textcolor{rust}{\sffamily\bfseries 3.}] Business indicator and
subcomponents, which form the primary part of calculations for operational risk
capital.
\end{enumerate}

Auditors and regulators require proof regarding controls and risk, which comes from
reports and documentation like committee meeting minutes.

In addition to Pillar 3 requirements, financial institutions often need to inform
regulators about breaches of conduct and significant operational risk events based
on specific criteria:

\noindent\begin{minipage}{\textwidth}
\renewcommand{\arraystretch}{1.35}
\begin{tabularx}{\textwidth}{@{}L{0.22\textwidth}Y@{}}
\rowcolor{navy} \thd{Criterion} & \thd{Events it captures}\\
{\tabfont \textbf{a) Reputation}} & {\tabfont Events that could have a significant impact on firm reputation.}\\
\rowcolor{paleblue}
{\tabfont \textbf{b) Resilience}} & {\tabfont Events that could negatively impact goods or services provided to customers, or cause them harm.}\\
{\tabfont \textbf{c) Materiality}} & {\tabfont Events that exceed loss or materiality thresholds.}\\
\rowcolor{paleblue}
{\tabfont \textbf{d) Stability}} & {\tabfont Events that could negatively impact the financial system or the ability to continue operations.}\\
\end{tabularx}
\end{minipage}
\orfigcap{The four criteria triggering regulatory notification.}
